% teilt dem Editor mit, dass er latin1 als encoding verwenden soll
% !TeX encoding = latin1

% festlegen und modifizieren der Dokumentenklasse
\documentclass[
	pdftex,						% PDFtex verwenden
	a4paper, 						% A4 Format benutzen
	12pt, 							% Schriftgr??e 12 f?r normalen Text
	headsepline						% Linie nach Kopfzeile
]{scrartcl}							% KOMA-Script Dokumentklasse "scrartcl" verwenden

% einbinden der Pakete

\usepackage{hyperref}
\usepackage{wrapfig}					% Text neben Bild
\usepackage{amsmath} 					% Paket f?r Formelumgebungen
\usepackage{amssymb} 					% Paket f?rMathe-Symbole

\usepackage{physics}                    %für Bra und Ket Notation

\usepackage{array}       					% Paket zum Erweitern der Tabelleneigenschaften
\usepackage{graphicx}  					% Paket um Grafiken einbetten zu k?nnen
\usepackage{color}       					% Paket f?r Farben im PDF (Seitenfarbe und Textfarbe)
\usepackage{makeidx}   					% Paket f?r die Indexerstellung.
\usepackage[T1]{fontenc} 
\usepackage[utf8]{inputenc}				% Paket f?r Schriftart, Verwenden von T1 Fonts in Ausgabe zur Darstellung von Umlauten
\usepackage{url}     					% Paket zur richtigen Darstellung von URLs
\usepackage[a4paper,left=2.5cm,right=2.0cm]{geometry} % Paket zum Einstellen der Seitenr?nder
\usepackage[ngerman]{babel} 				% Paket f?r Zeichensetzung und Worttrennung, deutsche Zeichensetzung verwenden
  					% Paket f?r Vektorpfeile: \vv{x}

%\usepackage{scrpage2} 
\usepackage{subfigure} 					%Bilder nebeneinander
\usepackage{caption} 					% Paket um Beschriftung f?r Gleitobjekte anzupassen
	\setlength{\captionmargin}{10pt} 		% extra Abstand rechts und links f?r Beschriftung
	\renewcommand{\captionfont}{\footnotesize}	% Beschriftung ist in Fu?notengr??e
			% Paket f?r erweiterte Positionierung von Gleitobjekten, Gleitobjekte d?rfen Abschnitt nicht verlassen
\usepackage{float}						% Paket zur erweiterung der Gleitobjektfunktionalit?ten, erm?glicht feste Positionierung mit \begin{table}[H]
\usepackage{epstopdf}			%Paket zur konvertierung von tif dateien zu png
\epstopdfDeclareGraphicsRule{.tif}{png}{.png}{convert #1 \OutputFile}
\AppendGraphicsExtensions{.tif}
% Setzen globaler Parameter
\graphicspath{{bilder/}} 					% legt globalen Pfad f?r Bilddateien fest
\setlength{\parindent}{0mm} 				% unterdr?ckt den Absatzeinzug


% passt Kopfzeile an
\pagestyle{myheadings}


%\usepackage{scrpage2} 
%\clearscrheadfoot
%\pagestyle{scrheadings}
%\automark[section]{section}
%\ihead{\rightmark}
%\ohead{\leftmark}
%\cfoot{\pagemark}


\begin{document}

% f?gt die Titelseite ein
\begin{titlepage}

\huge\textsc{Machine Learning 1}
\vspace{0,5cm}

\begin{minipage}{0.7\textwidth}
\large\textsc{Technical University Berlin\\ Winter semester 2020/21}
\end{minipage}
\begin{minipage}{0.3\textwidth}
\begin{flushright}
		\includegraphics[width=2cm]{pictures/tu-logo.jpg}
	\end{flushright}
\end{minipage}

\begin{center}

\vspace{0,5cm}
\hline

\vspace{2cm}
\Large\textbf{\textsc{}}
 \Large\textbf{\textsc{Exercise Sheet 1}}
\\
\vspace{0,7cm}
\small(09. November 2020)
\end{center}

\begin{table}[H]
\centering
\begin{tabular}{r l l}
								&								&			\\


\textbf{Students:} 	

								&	Mattes Ohse				 	& 337356        \\	
								&	Florian Ebert		 		& 391837 		\\
								&	Rodrigo Alexis Pardo Meza		 					&			    \\
								&   Bertty Contreras Rojas& \\



\end{tabular} 
\end{table}  
\vspace{3cm}
\normalsize
\begin{center}
   
\end{center}


% Inhaltsverzeichnis
\pagenumbering{Roman} 	% r?mische Zahlen (I,II,III,IV,...) f?r Inhaltsverzeichnis
\newpage

%\tableofcontents		% Inhaltsverzeichnis
%\listoffigures 		 	% Abbildungsverzeichnis

% der Rest ist arabisch (1,2,3,4,...) numeriert:
\pagebreak
\pagenumbering{arabic}	% aribische Numerierung
\setcounter{page}{1}	% setzt den Seitenz?hler auf "1"

% zusammenf?gen der Abschnitte
\section{Exercise 1: Maximum-Likelihood Estimation}
\section{Exercise 1: Estimating the Bayes Error}
\subsection{a}
The Bayes Formular is given by:
\begin{center}
    \begin{equation}
        P(w_{j})=\frac{P(x\mid w_{j})}{p(x)}.
    \label{eq:Bayesformula}
    \end{equation}
\end{center}

and the Bayes Error is given by:
\begin{center}
    \begin{equation}
        P(error)=min[P(w_{1}|x),P(w_{2}|x)].
        \label{eq:error}
    \end{equation}    
\end{center}

We now wanna show, that the full error can be upper-bounded as:

\begin{center}
    \begin{equation}
        P(error) &\leq\int\frac{2}{\frac{1}{P(w_{1}\mid x)}+\frac{1}{P(w_2|x)}} p(x) dx.
        \label{eq:1a}
    \end{equation}
\end{center}

We can assume that $P(w_{1}|x)\geq P(w_{2}|x)$, so that (\ref{eq:error}) becomes to $ P(error)=P(w_{2}|x)$. Now we can use (\ref{eq:1a}), under consideration that we are integrating on both sides for the same variable what means that we can ignore the integralsigns:

\begin{align*}
&\Rightarrow& P(error) &\leq\frac{2}{\frac{1}{P(w_{1}\mid x)}+\frac{1}{P(w_2|x)}} p(x)\\
&\Leftrightarrow& P(w_{2}|x)p(x)&\leq\frac{2}{\frac{1}{P(w_{1}\mid x)}+\frac{1}{P(w_2|x)}} p(x)\\
&\Leftrightarrow& P(w_{2}|x)&\leq\frac{2}{\frac{1}{P(w_{1}\mid x)}+\frac{1}{P(w_2|x)}}\\
&\Leftrightarrow& \frac{P(w_{2}|x)}{P(w_{1}|x)}+\frac{P(w_{2}|x)}{P(w_{2}|x)}&\leq2\\
&\Leftrightarrow& \frac{1}{P(w_{1}|x)}+\frac{1}{P(w_{2}|x)}&\leq\frac{2}{P(w_{2}|x)}\\
&\Leftrightarrow& \frac{1}{P(w_{1}|x)}&\leq\frac{1}{P(w_{2}|x)}\\
\\
&\Leftrightarrow& P(w_{1}|x)&\geq P(w_{2}|x).\\
\end{align*}
And that improves (\ref{eq:1a})






\subsection{b}
We want to show that for the distributions

\begin{center}
    \begin{equation}
        p(x|w_{1})=\frac{\pi^{-1}}{1+(x-\mu)^{2}},\text{ } p(x|w_{2})=\frac{\pi^{-1}}{1+(x+\mu)^{2}}
    \end{equation}
\end{center}
that the Bayes error can be upper-bounded by:

\begin{center}
\begin{equation}
    P(error)\leq\frac{2P(w_{1})P(w_{2})}{\sqrt{1+4\mu^{2}P(w_{1})P(w_{2})}}
\end{equation}
\end{center}

\begin{align*}
\Rightarrow P(error) &\leq\int\frac{2}{\frac{1}{P(w_{1}\mid x)}+\frac{1}{P(w_2|x)}} p(x) dx\\
&=\int\frac{2}{\frac{p(x)}{p(x\mid w_{1})\cdot P(w_{1})}+\frac{p(x)}{p(x|w_{2})\cdot P(w_{2})}} p(x) dx\\
&=\int\frac{2}{\frac{1}{p(x\mid w_{1})\cdot P(w_{1})}+\frac{1}{p(x|w_{2})\cdot P(w_{2})}} dx\\
&=\int\frac{2}{\frac{1}{\frac{\pi^{-1}}{1+(x-\mu)^{2}}P(w_{1})}+\frac{1}{\frac{\pi^{-1}}{1+(x+\mu)^{2}}P(w_{2})}}dx\\
&=\int\frac{2}{\frac{\pi((x-\mu)^{2}+1)}{P(w_{1})}+\frac{\pi((x+\mu)^{2}+1)}{P(w_{2})}}dx\\
&=\int\frac{2}{\frac{\pi(P(w_{2})((x-\mu)^{2}+1)+P(w_{1})((x+\mu)^{2}+1))}{P(w_{1})P(w_{2})}}\\
&=\int\frac{2\cdot P(w_{1})P(w_{2})}{\pi(P(w_{2})((x-\mu)^{2}+1)+P(w_{1})((x+\mu)^{2}+1))}\\
\\
&=\frac{2\cdot P(w_{1})P(w_{2})}{\pi}\int\frac{1}{P(w_{2})(x-\mu)^{2}+P(w_{2})+P(w_{1})(x+\mu)^{2}+P(w_{1})}dx\\
\\
&\text{Use } (x-\mu)^{2}=x^{2}-2x\mu+\mu^{2} \text{ and } (x+\mu)^{2}=x^{2}+2x\mu+\mu^{2}\\
\end{align*}
\\
\begin{align*}
=&\frac{2\cdot P(w_{1})P(w_{2})}{\pi}\\
&\cdot\int\frac{dx}{x^{2}P(w_{2})-2x\mu P(w_{2})+\mu^{2}P(w_{2})+P(w_{2})+x^{2}P(w_{1})+2x\mu P(w_{1})+\mu^{2}P(w_{1})+P(w_{1})}\\
\\
&=\frac{2\cdot P(w_{1})P(w_{2})}{\pi}\\
&\cdot\int\frac{dx}{x^{2}(P(w_{2})+P(w_{1}))+x2\mu(P(w_{1})-P(w_{2}))+\mu^{2}(P(w_{2})+P(w_{1}))+P(w_{2})+P(w_{1})}
\end{align*}

\begin{align*}
&\text{We say: }\\
&a=P(w_{1})+P(w_{2}), \\
&b=2\mu(P(w_{1}-P(w_{2}))\\ 
&c=\mu^{2}(P(w_{1})+P(w_{2}))+P(w_{1})+P(w_{2})\\
&\text{and we use: } \int\frac{1}{ax^2+bx+c}dx=\frac{2\pi}{\sqrt{4ac-b^{2}}}.\\
%&\text{Now we take a look at 4ac:}\\
\end{align*}
\begin{align*}
\text{Now we take a look at 4ac:}\\
4ac=&4[P(w_{1})+P(w_{2})]\cdot[\mu^{2}(P(w_{1})+P(w_{2}))+P(w_{1})+P(w_{2})]\\
=&4(\mu^{^2}+1)[P(w_{1})+P(w_{2})]^{2}.\\
\text{With that and the result of $b^{2}$}\\
b^{2}=&4\mu(P(w_{1})-P(w_{2}))^{2}= 4\mu^{2}[P^{2}(w_{1})-2P(w_{1})P(w_{2})+P^{2}(w_{2})]\\
\text{becomes the integral above to:}
\end{align*}

\begin{align*}
\Rightarrow&\frac{2P(w_{1})P(w_{2})}{\pi}\\
&\cdot\frac{2\pi}{\sqrt{4(\mu^{^2}+1)[P(w_{1})+P(w_{2})]^{2}-4\mu^{2}[P^{2}(w_{1})-2P(w_{1})P(w_{2})+P^{2}(w_{2})]}}\\
\\
=&\frac{2P(w_{1})P(w_{2})}{\sqrt{4\mu^{2}P(w_{1})P(w_{2})+2P(w_{1})P(w_{2})+P(w_{1})+P(w_{2})}}\vert(\text{use: $(a+b)^{2}=a^{2}+2ab+b^{2}$)}\\
\\
=&\frac{2P(w_{1})P(w_{2})}{\sqrt{4\mu^{2}P(w_{1})P(w_{2})+(P(w_{1})+P(w_{2}))^{2}}}\vert(\text{We know that: $P(w_{1})+P(w_{2})=1$})\\
\\
=&\frac{2P(w_{1})P(w_{2})}{\sqrt{1+4\mu^{2}P(w_{1})P(w_{2})}}
\end{align*}
This proves that the full error can be upper-bounded as:

\begin{center}
\begin{equation}
    P(error)\leq\frac{2P(w_{1})P(w_{2})}{\sqrt{1+4\mu^{2}P(w_{1})P(w_{2})}}
\end{equation}
\end{center}




\subsection{c}
For the two cases whether the data is low-dimensional or the data is high-dimensional we can find the upper-bounds by using the methods Chernoff Bound or Bhattacharyya Bound. Both of these methods are numerical. The functions of these integrals are usually not continuous, so that the bounds determined in this way are often tighter than the bounds determined analytically. In high-dimensional cases, Bhattacharyya's method should be better because it is less computationally expensive than the Chernoff Bound method.  
\newpage

\section{Exercise 2}
\begin{center}
    \begin{equation}
        D=(x_{1}, x_{2},..., x_{7})=(\textrm{head, head, tail, tail, head, head, head})
        \label{eq:sequence}
    \end{equation}
\end{center}

\begin{center}
        \[P(x|\theta)=\begin{cases}
        \theta & \text{if } x = \text{head}\\
        1-\theta& \text{if } x = \text{tail} \\
\end{cases}\]
\end{center}
where $\theta\in[0,1]$.

\subsection{a}
Following the lecture, the likelihood function $P(D|\theta)$ is given by:\\

\begin{align*}
    P(D|\theta)&=\prod_{k=1}^{N} P(x_{k}|\theta)\\
    &=\prod_{k=1}^{N_{\text{tail}}} P(\text{tail}|\theta)\cdot\prod_{k=1}^{N_{\text{Head}}} P(\text{head}|\theta)\\
    &=\prod_{k=1}^{2} (1-\theta)\cdot\prod_{k=1}^{5} \theta\\
    \\
    &=(1-\theta)^{2}\cdot\theta^{5}\\
\end{align*}

\subsection{b}
At first we determine the maximum likelihood solution $\hat{\theta}$. To find $\hat{\theta}$ we derive $P(D|\theta)$ (from exercise 2a) by $\theta$ and by setting $\frac{\partial}{\partial\theta}P(D|\theta)=0$ and solving this equation for $\theta$ we will find $\theta$.\\

\begin{align*}
    \frac{\partial}{\partial\theta} P(D|\theta) &= \frac{\partial}{\partial\theta}(1-\theta)^{2}\cdot\theta^{5}\\
    &= -2\cdot(1-\theta)\cdot\theta^{5}+5\cdot(1-\theta)^{2}\cdot\theta^{4}\\
    &= -12\cdot\theta^{5}+7\cdot\theta^{6}+5\cdot\theta^{4}\\
    &= \theta^{4}\cdot(7\cdot\theta^{2}-12\cdot\theta+5)
\end{align*}

\begin{align*}
    &\Rightarrow&&\frac{\partial}{\partial\theta} P(D|\theta)=0\\
    &\Leftrightarrow&&\theta^{4}\cdot(7\cdot\theta^{2}-12\cdot\theta+5)=0\\
\end{align*}
So we find the first theta: $\hat{\theta_{1}}=0$.\\

\begin{align*}
    &\Rightarrow&&7\cdot\theta^{2}-12\cdot\theta+5=0\\
    &\Leftrightarrow&&\theta^{2}-\frac{12}{7}\cdot\theta+\frac{5}{7}=0 \text{  }|\text{  p-q-formula}\\
    &\Leftrightarrow&&\hat{\theta}_{2,3}=\frac{12}{14}\pm\sqrt{\left(\frac{12}{14}\right)^{2}-\frac{5}{7}}\\
    &\Leftrightarrow&&\hat{\theta}_{2}=1, \text{  }\hat{\theta}_{3}=\frac{5}{7}\\
\end{align*}

So we have three posible solutions for $\hat{\theta}$:\\
\begin{center}
\begin{equation}
    \hat{\theta_{1}}=0\text{,  } \hat{\theta}_{2}=1\text{,  } \hat{\theta}_{3}=\frac{5}{7}.
\end{equation}
\end{center}
        
To decide which $\hat{\theta}$ is the maximum in the intervall $[0,1]$ we take a look at the functionvalues of $\frac{\partial^{2}}{\partial\theta^{2}}P(D|\hat{\theta}_{i})$.

\begin{align*}
    \frac{\partial^{2}}{\partial\theta^{2}}P(D|\hat{\theta}) &=\frac{\partial}{\partial\theta}\theta^{4}\cdot(7\cdot\theta^{2}-12\cdot\theta+5)\\
    &=4\cdot\theta^{3}\cdot(7\cdot\theta^{2}-12\cdot\theta+5)+\theta^{4}\cdot(14\cdot\theta-12)\\
    &=42\cdot\theta^{5}-60\cdot\theta^{4}+20\cdot\theta^{5}
\end{align*}

\begin{align*}
    &\Rightarrow\frac{\partial^{2}}{\partial\theta^{2}}P(D|\theta_{1}=0)=42\cdot\theta^{5}_{1}-60\cdot\theta^{4}_{1}+20\cdot\theta^{5}_{1}=0\rightarrow\text{saddle point}\\
    &\Rightarrow\frac{\partial^{2}}{\partial\theta^{2}}P(D|\theta_{2}=1)=42\cdot\theta^{5}_{2}-60\cdot\theta^{4}_{2}+20\cdot\theta^{5}_{2}=2\rightarrow\text{minima}\\
    &\Rightarrow\frac{\partial^{2}}{\partial\theta^{2}}P(D|\theta_{3}=\frac{5}{7})=42\cdot\theta^{5}_{3}-60\cdot\theta^{4}_{3}+20\cdot\theta^{5}_{3}\approx0,179\rightarrow\text{maxima}\\
\end{align*}

So the maximum likelihood solution for $\hat{\theta}$ is:
\begin{center}
    \begin{equation}
        \hat{\theta}=\frac{5}{7}
    \end{equation}
\end{center}

Now we can calculate the probability that the next two tosses are "head" under assuming that all tosses are generated independently:

\begin{align*}
    P(x_{8}=\text{head}, x_{8}=\text{head}|\theta)&=P(x_{8}=\text{head})\cdot P(x_{9}=\text{head})\\
    &=\hat{\theta}\cdot\hat{\theta}=\frac{25}{49}\approx0,51=51\protect
\end{align*}

\subsection{c}
Now the prior distribution for the parameter $\theta$ is definded as:\\
\begin{center}
        \[p(\theta)=\begin{cases}
        1 & \text{if } 0\le\theta\le1\\
        0& \text{else} \\
\end{cases}\]
\end{center}
    
\begin{align*}
    \Rightarrow p(\theta|D)&=\frac{p(D|\theta)p(\theta)}{\int p(D|\theta)p(\theta)d\theta}\text{      |}p(\theta)=1\\
    \\
    &=\frac{p(D|\theta)}{\int p(D|\theta)d\theta}\\
    \\
    &=\frac{(1-\theta)^{2}\cdot\theta^{5}}{\int_{0}^{1}(1-\theta)^{2}\cdot\theta^{5}d\theta}\\
    \\
    &=\frac{(1-\theta)^{2}\cdot\theta^{5}}{[\frac{\theta^{6}}{6}-\frac{2\cdot\theta^{7}}{7}+\frac{\theta^{8}}{8}]_{0}^{1}}\\
    \\
    &=168\cdot(1-\theta)^2\cdot\theta^{5}
\end{align*}

So we got the following form for $p(\theta|D)$:

\begin{center}
        \[p(\theta|D)=\begin{cases}
        168\cdot(1-\theta)^2\cdot\theta^{5}& \text{for } 0\le\theta\le1\\
        0& \text{else} \\
\end{cases}\]
\end{center}

With this result for $p(\theta|D)$ we can now calculate the probability that the next two tosses are head:

\begin{align*}
    &\Rightarrow\int P(x_{8}=\text{head},x_{9}=\text{head}|\theta)\cdot p(\theta|D) d\theta\\
    &=\int \underbrace{P(x_{8}=\text{head}|\theta)}_{=\theta}\cdot \underbrace{P(x_{9}=\text{head}|\theta)}_{=\theta}\cdot \underbrace{p(\theta|D)}_{=168\cdot(1-\theta)^{2}\cdot\theta^{5}} d\theta\\
    &=168\cdot\int_{0}^{1}\theta^{7}(1-\theta)^{2}d\theta\\
    &=168\cdot\int_{0}^{1}(\theta^{9}-2\cdot\theta^{8}+\theta^{7})d\theta\\
    &=168\cdot\left[\frac{\theta^{10}}{10}-\frac{2\cdot\theta^{9}}{9}+\frac{\theta^{8}}{8}\right]_{0}^{1}\\
    &=\frac{7}{15}\approx0,467=46,7\protect
\end{align*}
\section{Exercise 3}
Given formulas are:\\
\begin{center}
    \begin{equation}
        \frac{1}{\sigma^{2}_{n}}=\frac{n}{\sigma^{2}}+\frac{1}{\sigma^{2}_{0}}
        \label{eq:sigma}
    \end{equation}
\end{center}
\begin{center}
    \begin{equation}
        \frac{\mu_{n}}{\sigma^{2}_{n}}=\frac{n}{\sigma^{2}_{n}}\hat{\mu}_{n}+\frac{\mu_{0}}{\sigma^{2}_{0}}
    \end{equation}
\end{center}
\begin{center}
    \begin{equation}
        \hat{\mu}_{n}=\frac{1}{n}\sum_{k=1}^{n}x_{k}
    \end{equation}
\end{center}


\subsection{a}
We want to show that following relation holds:

\begin{center}
    \begin{equation}
        \sigma^{2}_{n}\le min\left(\frac{\sigma^{2}}{n},\sigma^{2}_{0}\right)
    \label{eq:3a}
    \end{equation}
\end{center}

At first we need a expression for $\sigma^{2}_{n}$. This will give us formula (\ref{eq:sigma}):

\begin{align*}
    &\Rightarrow&\frac{1}{\sigma^{2}_{n}} &=\frac{n}{\sigma^{2}}+\frac{1}{\sigma^{2}_{0}}\\
    &\Leftrightarrow& 1 &=\left(\frac{n}{\sigma^{2}}+\frac{1}{\sigma^{2}_{0}}\right)\cdot\sigma^{2}_{n}\\
    &\Leftrightarrow& \sigma^{2}_{n} &=\frac{1}{\frac{n}{\sigma^{2}}+\frac{1}{\sigma^{2}_{0}}}\\
    &\Leftrightarrow& \sigma^{2}_{n} &=\frac{1}{\frac{n\sigma^{^2}_{0}}{\sigma^{2}\sigma^{^2}_{0}}+\frac{\sigma^{2}}{\sigma^{2}_{0}\sigma^{2}}}\\
    &\Leftrightarrow& \sigma^{2}_{n} &=\frac{1}{\frac{n\sigma^{^2}_{0}+\sigma^{2}}{\sigma^2\sigma^{2}_{0}}}\\
    &\Leftrightarrow& \sigma^{2}_{n} &=\frac{\sigma^2\sigma^{2}_{0}}{n\sigma^{^2}_{0}+\sigma^{2}}
\end{align*}

Now we have two cases to consider to show that $\sigma^{2}_{n}\le min\left(\frac{\sigma^{2}}{n},\sigma^{2}_{0}\right)$ holds:
\begin{align}
    \text{case 1: } \sigma^{2}_{0}&\le\frac{\sigma^{2}}{n}\\
    \text{case 2: } \sigma^{2}_{0}&\ge\frac{\sigma^{2}}{n}
\end{align}

At first we take a look at case one:\\

\begin{align*}
    &\Rightarrow&\sigma^{2}_{n}&\le min\left(\frac{\sigma^{2}}{n},\sigma^{2}_{0}\right)\\
    &\Leftrightarrow&\sigma^{2}_{n}&\le\frac{\sigma^{2}}{n}\\
    &\Leftrightarrow& \frac{\sigma^2\sigma^{2}_{0}}{n\sigma^{^2}_{0}+\sigma^{2}} &\le\frac{\sigma^{2}}{n}\\
    &\Leftrightarrow& \frac{\sigma^{2}_{0}}{n\sigma^{^2}_{0}+\sigma^{2}} &\le\frac{1}{n}\\
    &\Leftrightarrow& \frac{1}{n\sigma^{^2}_{0}+\sigma^{2}} &\le\frac{1}{n\sigma^{2}_{0}}\\
    &\Leftrightarrow& 1 &\le\frac{n\sigma^{^2}_{0}+\sigma^{2}}{n\sigma^{2}_{0}}\\
    &\Leftrightarrow& 1 &\le 1+\frac{\sigma}{n\sigma^{2}_{0}}\\
    &\Leftrightarrow& 0 &\le \frac{\sigma}{n\sigma^{2}_{0}}
\end{align*}

This expression is true because $\sigma$, n and $\sigma_{0}$ are always greater than zero (or maybe sometimes equal to zero) and so we have shown thas the first case is true.\\
Now we take a look at case two:

\begin{align*}
    &\Rightarrow&\sigma^{2}_{n}&\le min\left(\frac{\sigma^{2}}{n},\sigma^{2}_{0}\right)\\
    &\Leftrightarrow&\sigma^{2}_{n}&\le\sigma^{2}_{0}\\
    &\Leftrightarrow& \frac{\sigma^2\sigma^{2}_{0}}{n\sigma^{^2}_{0}+\sigma^{2}} &\le\sigma^{2}_{0}\\
    &\Leftrightarrow& \frac{\sigma^2}{n\sigma^{^2}_{0}+\sigma^{2}} &\le 1\\
    &\Leftrightarrow&\sigma^2&\le n\sigma^{^2}_{0}+\sigma^{2}\\
    &\Leftrightarrow& 1\le\frac{n\sigma^{2}_{0}}{\sigma^{2}}-1\\
    &\Leftrightarrow& 0\le\frac{n\sigma^{2}_{0}}{\sigma^{2}}
\end{align*}

Like in case one is this expression true because $\sigma$, n and $\sigma_{0}$ are always greater than zero (or maybe sometimes equal to zero) and so we have shown thas the second case is true.\\
So we have shown that relation (\ref{eq:3a}) holds.

\subsection*{3b}
Prove that $min(\hat{\mu}_n, \mu_0)\leq \mu_n \leq max(\hat{\mu}_n, \mu_0)$
\begin{align*}
   \frac{\mu_n}{\sigma^{2}_{n}} &= \frac{n}{\sigma^{2}}\hat{\mu}_n+\frac{\mu_0}{\sigma^{2}_0}\\
   \Leftrightarrow\frac{\mu_n}{\frac{\sigma^2\sigma^{2}_{0}}{n\sigma^{^2}_{0}+\sigma^{2}}} &= \frac{n}{\sigma^{2}}\hat{\mu}_n+\frac{\mu_0}{\sigma^{2}_0}\\
   \Leftrightarrow \mu_n&=\frac{n\sigma^{2}_0\hat{\mu}_n + \mu_0\sigma^{^2}}{n\sigma^{^2}_{0}+\sigma^{2}}
\end{align*}
\hspace*{2mm}If $\hat{\mu}_n \leq \mu_n$:
\begin{align*}
\hat{\mu}_n &\leq\frac{n\sigma^{2}_0\hat{\mu}_n + \mu_0\sigma^{^2}}{n\sigma^{^2}_{0}+\sigma^{2}}\\
\Leftrightarrow0&\leq\frac{n\sigma^{2}_0\hat{\mu}_n + \mu_0\sigma^{^2}-\hat{\mu}_n(n\sigma^{^2}_{0}+\sigma^{2})}{n\sigma^{^2}_{0}+\sigma^{2}}\\
\Leftrightarrow0&\leq\frac{\mu_0\sigma^{2}-\hat{\mu}_n\sigma^{2}}{n\sigma^{^2}_{0}+\sigma^{2}}\\
\Leftrightarrow 0&\leq \sigma^{2}(\mu_0-\hat{\mu}_n),\ because\ neither\ n,\ \sigma^{^2}_{0}\ or\ \sigma^{2}\ can\ be\ negative\\
\Leftrightarrow 0&\leq \mu_0-\hat{\mu}_n,\ because\ \sigma^{2}>0\\
\Leftrightarrow \hat{\mu}_n&\leq \mu_0
\end{align*}
\hspace{10mm}If $\hat{\mu}_n \leq \mu_n ,\ \mu_0$ has to be greater or equals to $\mu_n$ to fulfill \\\hspace*{10mm}$min(\hat{\mu}_n, \mu_0)\leq \mu_n \leq max(\hat{\mu}_n, \mu_0)$ because of transitivity\\ 
\hspace*{10mm}that would mean $\hat{\mu}_n\leq \mu_0$ is true.\\ \\
\hspace*{5mm}Test whether $\mu_0\geq \mu_n$ if $\hat{\mu}_n \leq \mu_n$:
\begin{align*}
    \mu_0&\geq \mu_n \\
    \Leftrightarrow \mu_0&\geq\frac{n\sigma^{2}_0\hat{\mu}_n + \mu_0\sigma^{^2}}{n\sigma^{^2}_{0}+\sigma^{2}}\\
    \Leftrightarrow 0&\geq\frac{n\sigma^{2}_0\hat{\mu}_n + \mu_0\sigma^{^2}-\mu_0(n\sigma^{^2}_{0}+\sigma^{2})}{n\sigma^{^2}_{0}+\sigma^{2}}\\
    \Leftrightarrow 0&\geq\frac{n\sigma^{2}_0\hat{\mu}_n -\mu_0n\sigma^{^2}_{0}}{n\sigma^{2}_{0}+\sigma^{2}}\\
    \Leftrightarrow 0&\geq n\sigma^{2}_0\hat{\mu}_n -\mu_0n\sigma^{2}_{0},\ because\ n,\ \sigma^{2}_{0}\ and\  \sigma^{2}>0\\
    \Leftrightarrow \mu_0&\geq \hat{\mu}_n
\end{align*}
\hspace*{5mm}$\Rightarrow$ If $\mu_0\geq \hat{\mu}_n$ then $min(\hat{\mu}_n, \mu_0)\leq \mu_n \leq max(\hat{\mu}_n, \mu_0)$ is true.\\ \\
\hspace*{2mm} Same can be proven for $\hat{\mu}_n \geq \mu_n$:
\begin{align*}
\hat{\mu}_n &\geq\frac{n\sigma^{2}_0\hat{\mu}_n + \mu_0\sigma^{^2}}{n\sigma^{^2}_{0}+\sigma^{2}}\\
\Leftrightarrow0&\geq\frac{n\sigma^{2}_0\hat{\mu}_n + \mu_0\sigma^{^2}-\hat{\mu}_n(n\sigma^{^2}_{0}+\sigma^{2})}{n\sigma^{^2}_{0}+\sigma^{2}}\\
\Leftrightarrow0&\geq\frac{\mu_0\sigma^{2}-\hat{\mu}_n\sigma^{2}}{n\sigma^{^2}_{0}+\sigma^{2}}\\
\Leftrightarrow 0&\geq \sigma^{2}(\mu_0-\hat{\mu}_n),\ because\ neither\ n,\ \sigma^{^2}_{0}\ or\ \sigma^{2}\ can\ be\ negative\\
\Leftrightarrow 0&\geq \mu_0-\hat{\mu}_n,\ because\ \sigma^{2}>0\\
\Leftrightarrow \hat{\mu}_n&\geq \mu_0
\end{align*}
\hspace*{5mm}Test whether $\mu_n\geq \mu_0$ if $\hat{\mu}_n \geq \mu_n$:
\begin{align*}
    \mu_0&\leq \mu_n \\
    \Leftrightarrow \mu_0&\leq\frac{n\sigma^{2}_0\hat{\mu}_n + \mu_0\sigma^{^2}}{n\sigma^{^2}_{0}+\sigma^{2}}\\
    \Leftrightarrow 0&\leq\frac{n\sigma^{2}_0\hat{\mu}_n + \mu_0\sigma^{^2}-\mu_0(n\sigma^{^2}_{0}+\sigma^{2})}{n\sigma^{^2}_{0}+\sigma^{2}}\\
    \Leftrightarrow 0&\leq\frac{n\sigma^{2}_0\hat{\mu}_n -\mu_0n\sigma^{^2}_{0}}{n\sigma^{2}_{0}+\sigma^{2}}\\
    \Leftrightarrow 0&\leq n\sigma^{2}_0\hat{\mu}_n -\mu_0n\sigma^{2}_{0},\ because\ n,\ \sigma^{2}_{0}\ and\  \sigma^{2}>0\\
    \Leftrightarrow \mu_0&\leq \hat{\mu}_n
\end{align*}
\hspace*{5mm}$\Rightarrow$ If $\hat{\mu}_n\geq \mu_n$ then $min(\hat{\mu}_n, \mu_0)\leq \mu_n \leq max(\hat{\mu}_n, \mu_0)$ is true.\\ \\
It is proven that $min(\hat{\mu}_n, \mu_0)\leq \mu_n \leq max(\hat{\mu}_n, \mu_0)$

\end{document}